% Modelos lineales por tramos del diodo
\begin{tikzpicture}[font=\footnotesize,
  ejes/.style={-{Stealth[length=4pt]},grisazul},
  real/.style={grisclaro,line width=2.2pt},
  modelo/.style={azul,very thick}]
  \newcommand{\ejesmodelo}[2]{%
    \draw[ejes] (-1.6,0) -- (2.3,0) node[below left]{$\VD$};
    \draw[ejes] (0,-0.4) -- (0,2.6) node[left]{$\ID$};
    \draw[real,domain=-1.6:1.33,samples=60] plot (\x,{0.0006*exp(\x/0.16)});
    \node[font=\bfseries] at (-1.4,2.5) {#1};
    \node[navy] at (0.4,-0.75) {#2};}
  \begin{scope}
    \ejesmodelo{(a)}{Ideal}
    \draw[modelo] (-1.6,0) -- (0,0) -- (0,2.4);
  \end{scope}
  \begin{scope}[xshift=4.6cm]
    \ejesmodelo{(b)}{Tensión umbral}
    \draw[modelo] (-1.6,0) -- (1.2,0) -- (1.2,2.4);
    \draw[grisazul,densely dotted] (1.2,0) -- (1.2,-0.15) node[below]{$\Vg$};
  \end{scope}
  \begin{scope}[xshift=9.2cm]
    \ejesmodelo{(c)}{Umbral $+\,r_d$}
    \draw[modelo] (-1.6,0) -- (1.1,0) -- (1.5,2.4);
    \draw[grisazul,densely dotted] (1.1,0) -- (1.1,-0.15) node[below]{$\Vg$};
    \draw[grisazul] (1.3,1.2) -- (1.7,1.2) node[right]{pendiente $1/r_d$};
  \end{scope}
\end{tikzpicture}
